\subsection{Cyclic extensions}

DEFINITION 3. --- An extension E of K is said to be cyclic if it is Galois and its Galois group is cyclic.

Examples. --- 1) Every Galois extension of prime degree is cyclic, because every finite group G of prime order p is cyclic (for every element \(x \neq 1\) of G is of order p, hence generates G).

\begin{enumerate}
\def\labelenumi{\arabic{enumi})}
\setcounter{enumi}{1}
\item
  Let \(\mathrm{F}(\mathbf{X}) = \mathbf{X}^2 + a\mathbf{X} + b\) be an irreducible polynomial in \(\mathbf{K}[\mathbf{X}]\) . The only case where \(\mathrm{F}(\mathbf{X})\) is not separable is that where \(\mathbf{K}\) is of characteristic 2 and \(a = 0\) . We shall leave this case aside; let \(\mathbf{E}\) be an extension of \(\mathbf{K}\) generated by a root \(x\) of \(\mathrm{F}(\mathbf{X})\) . We have \([\mathrm{E}:\mathrm{K}] = 2\) and \(\mathrm{F}(\mathbf{X}) = (\mathbf{X} - x)(\mathbf{X} + a + x)\) , hence \(\mathbf{E}\) is a Galois extension of \(\mathbf{K}\) . The Galois group of \(\mathbf{E}\) over \(\mathbf{K}\) is of order 2, hence cyclic.
\item
  Let F be a field and \(\sigma\) an automorphism of finite order. The field E of invariants of \(\sigma\) is also the field of invariants of the cyclic group of order n generated by \(\sigma\) , and hence (V, p.~66, Th. 3) F is a cyclic extension of degree n of E.
\end{enumerate}

We know (I, p.~50) that every subgroup and every quotient group of a cyclic group is cyclic. Therefore (V, p.~68, Cor. 4), if E is a cyclic extension of a field K, of degree n, then every subextension F of E is cyclic over K, and E is cyclic over F. For every divisor d of n there exists a unique subfield F of degree d over K contained in E: for in a cyclic group of order n there exists a unique subgroup of index d.

THEOREM 3 (Hilbert). --- Let E be a cyclic extension of K and let \(\sigma\) be a generator of the Galois group \(\Gamma\) of E over K.

\begin{enumerate}
\def\labelenumi{\alph{enumi})}
\item
  For an element \(x \in \mathbb{E}\) to satisfy \(\mathrm{N}_{\mathrm{E} / \mathrm{K}}(x) = 1\) it is necessary and sufficient that \(y \in \mathrm{E}^*\) exist such that \(x = y / \sigma(y)\) ; every \(y_1 \in \mathrm{E}^*\) with \(x = y_1 / \sigma(y_1)\) is then of the form \(\lambda y\) where \(\lambda \in \mathbf{K}^*\) .
\item
  For an element \(x \in E\) to satisfy \(\operatorname{Tr}_{\operatorname{E}/\operatorname{K}}(x) = 0\) it is necessary and sufficient that \(z \in E\) exist with \(x = z - \sigma(z)\) ; then every \(z_{1} \in E\) with \(x = z_{1} - \sigma(z_{1})\) is of the form \(z + \mu\) , where \(\mu \in K\) .
\end{enumerate}

Let us first prove a lemma.

Lemma 4. --- Let \(\Gamma\) be a cyclic group of order \(n\) , \(\sigma\) a generator of \(\Gamma\) and M a commutative group on which \(\Gamma\) operates according to the rule \(\gamma\) . \((m + m') = \gamma\) . \(m + \gamma\) . \(m'\) for all \(\gamma \in \Gamma\) and \(m\) , \(m' \in M\) . Let Z be the set of all mappings of \(\Gamma\) into M satisfying the relation

\[
f \left(\tau \tau^ {\prime}\right) = f (\tau) + \tau . f \left(\tau^ {\prime}\right) \quad \text { for } \quad \tau , \tau^ {\prime} \text { in } \Gamma .\tag{10}
\]

Put \(u(f) = f(\sigma)\) for \(f \in \mathbf{Z}\) and \(t(m) = \sum_{\tau \in \Gamma} \tau \cdot m\) for \(m \in \mathbf{M}\) . Then the sequence

\[
0 \to \mathbf {Z} \xrightarrow {u} \mathbf {M} \xrightarrow {t} \mathbf {M}
\]

is exact.

Let \(f \in Z\) ; taking \(\tau = \tau' = 1\) in (10) we find that \(f(1) = 0\) . Further, by induction on \(m \geqslant 0\) we deduce the relation

\[
f (\sigma^ {m}) = f (\sigma) + \sigma \cdot f (\sigma) + \dots + \sigma^ {m - 2} \cdot f (\sigma) + \sigma^ {m - 1} \cdot f (\sigma).\tag{11}
\]

We have \(\sigma^n = 1\) , whence \(f(\sigma^n) = 0\) ; the preceding relation with \(m = n\) is equivalent to the equality \(t(u(f)) = 0\) , whence \(\operatorname{Im} u \subset \ker t\) . Moreover, from (11) it follows that \(\operatorname{Ker} u = 0\) .

Let \(m \in M\) be such that \(t(m) = 0\) , that is, \(m + \sigma \cdot m + \cdots + \sigma^{n-1} \cdot m = 0\) . Let us define the mapping \(f\) of \(\Gamma\) into \(M\) by

\[
f \left(\sigma^ {j}\right) = m + \sigma . m + \dots + \sigma^ {j - 2}. m + \sigma^ {j - 1}. m\tag{12}
\]

for \(0 \leqslant j \leqslant n - 1\) . It may be left to the reader to establish the relation (10). We clearly have \(m = f(\sigma)\) , whence \(\operatorname{Im} u \supset \operatorname{Ker} t\) .

With the lemma now proved, let us take \(y \in E^{*}\) and \(x = y / \sigma(y)\) ; we have \(\mathrm{N}_{\mathrm{E}/\mathrm{K}}(\sigma(y)) = \mathrm{N}_{\mathrm{E}/\mathrm{K}}(y)\) , whence \(\mathrm{N}_{\mathrm{E}/\mathrm{K}}(x) = 1\) . Conversely, let x in \(E^{*}\) be such that \(\mathrm{N}_{\mathrm{E}/\mathrm{K}}(x) = 1\) ; by Lemma 4, applied to \(M = E^{*}\) , there exists a family of elements \((c_{\tau})_{\tau \in \Gamma}\) of \(E^{*}\) satisfying the relation \(c_{\tau\tau'} = c_{\tau} \cdot \tau(c_{\tau'})\) for \(\tau\) , \(\tau'\) in \(\Gamma\) and \(c_{\sigma} = x\) . By Cor. 1 of Prop. 9 (V, p.~65) there exists \(y \in E^{*}\) with \(c_{\tau} = y / \tau(y)\) for all \(\tau \in \Gamma\) , whence in particular \(x = c_{\sigma} = y / \sigma(y)\) . If \(y_{1} \in E^{*}\) satisfies \(x = y_{1} / \sigma(y_{1})\) , then we have

\[
\sigma \left(y _ {1} y ^ {- 1}\right) = y _ {1} y ^ {- 1},
\]

hence \(y_{1}y^{-1}\) belongs to \(K^{*}\) because \(\sigma\) generates the Galois group of E over K. This proves a).

The assertion b) follows in similar fashion from Cor. 2 of Prop. 9 (V, p.~65).

